% \section{Introduction}

% The era of diffusion models~\cite{song2020score, song2020improved, lipman2022flow} has heralded a period of monumental development in video generation. The trajectory began with foundational UNet-based diffusion models~\cite{guo2023animatediff,singer2022make,ho2022video,xing2024dynamicrafter} and has recently culminated in DiT-based architectures~\cite{kong2024hunyuanvideo,yang2024cogvideox,wan2025wan,polyak2024movie,ma2025step}, pioneered by landmark models like Sora~\cite{openaisora2024}. This new generation of video models has achieved unprecedented progress in both prompt alignment and generation quality. These exhilarating advancements are catalyzing a profound transformation in the field of video creation, unlocking limitless possibilities.
% Commercial applications, such as film production, advertising, and creative short-form video generation, stand to benefit immensely. By leveraging video generation technology, professionals can significantly enhance their workflow efficiency and reduce production costs. Concurrently, the evolution of these powerful video models has spurred a wave of research into downstream applications built upon pre-trained backbones. Fields such as talking head synthesis~\cite{wang2023progressive,yu2023talking,wang2025fantasytalking,luo2025dreamactor}, video-based virtual try-on~\cite{ning2024picture,wang2025mv,kim2024stableviton}, human animation~\cite{tan2024animate,chen2025hunyuanvideo,lin2025omnihuman}, and video editing~\cite{xue2025stand,jiang2025vace} have all seen remarkable breakthroughs, empowered by the enhanced capabilities of modern video models. These applications are now providing viable solutions for a multitude of real-world scenarios.

% The current generation of video synthesis, driven by diffusion foundation models, has evolved along several key axes. A primary development has been in model architecture, marked by a transition from early UNet-based frameworks~\cite{rombach2022high} to the now-prevalent Diffusion Transformer (DiT) architecture~\cite{peebles2023scalable}. Initial models, such as AnimateDiff~\cite{guo2023animatediff} and Make-A-Video~\cite{singer2022make}, extended pre-trained UNet image generators by incorporating temporal layers to enforce inter-frame consistency. A paradigm shift occurred with the advent of Sora~\cite{openaisora2024}, followed by open-source DiT models like CogVideox~\cite{yang2024cogvideox}, HunyuanVideo~\cite{kong2024hunyuanvideo}, and WAN2.1~\cite{wan2025wan}. These models, trained from scratch with 3D full-attention mechanisms, have demonstrated significantly more promising results.

% Furthermore, the success of Sora has underscored the importance of scaling laws for both data and model size. A key insight for the research community is that systematically scaling up a DiT-based architecture leads to substantial improvements in video generation performance. This principle has been validated not only by open-source efforts but also by a series of high-performing, closed-source models such as SeeDance1.0~\cite{gao2025seedance}, Kling~\cite{kuaishou2024kling}, Veo2~\cite{veo22025}, etc., all of which leverage the DiT architecture at a massive scale.

% Architecturally, these models share a common blueprint: a 3D Variational Autoencoder (VAE) first compresses the video into a latent representation. The DiT then operates on this latent space, executing a conditional denoising process on noisy latents to reverse the diffusion and generate the final video. The development efforts for these models have predominantly focused on two objectives: enhancing semantic alignment to text prompts and improving the overall visual quality of the output. To this end, researchers have employed a suite of techniques, including curating high-quality datasets, optimizing video captioning models, advanced prompt engineering, integrating super-resolution modules, and implementing post-training alignment strategies such as Supervised Fine-Tuning (SFT)~\cite{instag} and Reinforcement Learning from Human Feedback (RLHF)~\cite{rlhf}. These validated methods have continuously advanced model capabilities and, in turn, have empowered content creators by streamlining the initial video production workflow, thereby reducing time and cost while boosting creative efficiency.

% Despite these advancements, a significant limitation persists: a near-exclusive focus on the visual modality. Video is an inherently multimodal medium, combining both visual and audio streams. While post-hoc video-to-audio models can generate sound for synthesized videos, this decoupled approach is fundamentally flawed. It can adapt the audio to the video content, but it cannot enforce temporal alignment in the reverse direction. For example, it is incapable of synchronizing the lip movements of a speaking character with the corresponding speech. Google's Veo3~\cite{veo32025} has demonstrated the ability to generate audio and video synchronously, but it remains a closed-source model with no publicly available technical details.

% To address this critical gap between closed-source commercial systems and open research, we introduce \textbf{UniVerse-1}: a unified, fully open-source, Veo-3-like model capable of simultaneously generating coordinated audio and video. 

% To enhance training efficiency and capitalize on the robust capabilities of existing pre-trained models, our approach circumvents training a new model from scratch. Instead, we propose a \textbf{stitching of experts (SoE)} methodology that integrates a state-of-the-art video generation model, WAN2.1~\cite{wan2025wan}, with a music generation model, Ace-step~\cite{gong2025ace}. The integration of two experts is realized by introducing cross-modal MLP connectors within each corresponding block of the two models, enabling bidirectional interaction between the video and audio modalities. This strategy was found to significantly accelerate training convergence.
% To preserve the critical temporal and semantic alignment between audio-video data and their textual descriptions, we developed an \textbf{online annotation pipeline}. This pipeline generates labels dynamically during the training process, thereby mitigating the performance degradation that can arise from static, misaligned annotations.
% Furthermore, our investigation revealed a crucial insight for bimodal diffusion modeling: the necessity of independent noise sampling for each modality. We identified that without such isolation, the pseudo-random number generation process~\cite{hamming1952mathematical} can introduce spurious correlations between the noise vectors for video and audio, which subsequently degrades audio quality during inference. To support this work, we curated a high-quality dataset comprising approximately $7,600$ hours of precisely aligned audio-video content. To systematically evaluate our proposed method, we introduce \textbf{Verse-Bench}, a new benchmark dataset. Verse-Bench comprises $600$ image-text prompt pairs designed to cover a diverse range of sound categories. Highlighting its versatility, the benchmark supports the evaluation of not only joint audio-video generation but also unidirectional tasks. It can be used to assess video-to-audio generation, and it includes the specialized Verse-Ted subset for evaluating audio-to-video tasks like talking head synthesis.

% In summary, our primary contributions are:
% \begin{itemize}[leftmargin=*]
%     \item An Open-source Audio-Video Foundation Model: We present a novel, open-source model for synchronous audio-video generation that produces highly coherent and well-aligned audio-visual content.
%     \item A Pre-trained Model Fusion Strategy: We introduce a stitching of experts technique that effectively fuses the capabilities of existing pre-trained models, significantly accelerating convergence while leveraging their powerful priors.
%     \item An Online Data Annotation Pipeline: We developed an online data annotation pipeline that ensures strict temporal and semantic alignment between training data and captions, a crucial factor for high-fidelity bimodal generation.
%     \item A new Verse-Bench is proposed to comprehensively evaluate joint audio-video generation models.
% \end{itemize}


\section{Introduction}

The era of diffusion models~\cite{song2020score, song2020improved, lipman2022flow} has culminated in the rise of Diffusion Transformer (DiT) architectures~\cite{peebles2023scalable}, exemplified by landmark models like Sora~\cite{openaisora2024} and its open-source counterparts~\cite{kong2024hunyuanvideo,yang2024cogvideox,wan2025wan}. These models, leveraging unprecedented scales of data and computation, have achieved remarkable quality and prompt alignment in video generation. This success is catalyzing a profound transformation across creative industries and has spurred a wave of research into downstream applications such as talking head synthesis~\cite{wang2023progressive,yu2023talking,wang2025fantasytalking,luo2025dreamactor} and human animation~\cite{tan2024animate,chen2025hunyuanvideo,lin2025omnihuman}, which now offer viable real-world solutions.

However, this rapid progress has been almost exclusively confined to the visual domain, treating video as a silent movie. This \textbf{unimodal focus} represents a fundamental bottleneck, as it ignores the inherently multimodal nature of video. Post-hoc video-to-audio models~\cite{shan2025hunyuanvideo} serve as a superficial fix, but they are inherently limited; while capable of adapting audio to existing visual content, they fail to enforce temporal alignment in the reverse direction. This makes critical tasks, such as synchronizing lip movements with speech, impossible. While closed-source systems like Google's Veo3~\cite{veo32025} have demonstrated synchronous audio-video generation, the lack of publicly available technical details leaves a critical gap in open research.

To bridge this gap between closed-source systems and open research, we introduce \textbf{UniVerse-1}: a unified, fully open-source, Veo-3-like model capable of simultaneously generating coordinated audio and video. Our work is underpinned by several key technical contributions designed to address the unique challenges of bimodal generation.

Instead of the costly process of training a new model from scratch, we propose a novel and efficient \textbf{stitching of experts (SoE)} paradigm. This methodology effectively fuses a state-of-the-art video generation model, WAN2.1~\cite{wan2025wan}, with a music generation model, Ace-step~\cite{gong2025ace}. The core of this fusion lies in lightweight, cross-modal MLP connectors introduced within corresponding blocks of each model. These connectors facilitate bidirectional interaction between modalities, and we found this strategy to significantly accelerate training convergence by leveraging the powerful priors of the pre-trained experts.

Furthermore, we tackle the critical challenge of data alignment in bimodal training. We argue that static, pre-processed annotations are a flawed paradigm for tasks requiring precise temporal consistency. To address this, we developed an \textbf{online annotation pipeline} that generates labels dynamically during training. This approach ensures strict temporal and semantic alignment between audio-video data and their textual descriptions, mitigating the performance degradation caused by static misalignment. During our investigation, we also uncovered a crucial, yet overlooked, factor in bimodal diffusion modeling: \textbf{cross-modal noise correlation}. We identified that the standard pseudo-random number generation process~\cite{hamming1952mathematical} can introduce spurious correlations between the noise vectors for video and audio, which subsequently degrades audio quality during inference. Our solution involves ensuring independent noise sampling for each modality.

To support this work, we curated a high-quality dataset comprising approximately $7,600$ hours of precisely aligned audio-video content. To systematically evaluate our method, we also propose \textbf{Verse-Bench}, a new benchmark featuring $600$ image-text prompt pairs covering a diverse range of sound categories. Highlighting its versatility, Verse-Bench supports not only joint audio-video generation but also unidirectional tasks, including a specialized \textbf{Verse-Ted} subset designed for evaluating audio-to-video synthesis.

In summary, our primary contributions are:
\begin{itemize}[leftmargin=*]
    \item \textbf{An Open-source Audio-Video Foundation Model:} We present UniVerse-1, a novel, open-source model capable of producing highly coherent and well-aligned synchronous audio-visual content, closing a critical gap in the open-source community.
    
    \item \textbf{A Novel Methodology for Joint Audio-Video Generation:} We propose a comprehensive methodology to enable efficient and high-quality joint audio-video synthesis. This is achieved through three key innovations: a \textbf{stitching of experts (SoE)} paradigm to accelerate convergence by fusing pre-trained models; an \textbf{online data annotation pipeline} to solve the critical static misalignment problem in training; and the identification and mitigation of the previously overlooked \textbf{cross-modal noise correlation} issue, a crucial factor for generation quality.
    
    \item \textbf{A Comprehensive Evaluation Benchmark:} We propose Verse-Bench, a new benchmark designed to comprehensively evaluate joint audio-video generation models across a diverse set of tasks.
\end{itemize}